\chapter{Prediction, coordinate support, and canonical metrics}
\label{ch:prediction}
\chapterpromise{Family-specific posterior prediction, explicit extrapolation, and evaluation rules that do not create false success from unmatched boundaries.}

\section{Posterior prediction}
A prediction path must dispatch to the fitted observation family's posterior predictive. It may not
fall through to a Gaussian default. The exported model also carries a coordinate support interval.
Coordinates outside that interval raise an error by default; a named extrapolation policy may be
used only when its behavior and provenance are recorded.

\begin{figure}[H]
\centering
\resizebox{0.96\textwidth}{!}{\begin{tikzpicture}[node distance=7mm and 8mm]
\node[secondarybox,text width=29mm] (post) {Posterior distribution for each fitted segment};
\node[primarybox,right=of post,text width=29mm] (cond) {Condition on a MAP partition or average over partitions};
\node[primarybox,right=of cond,text width=29mm] (coord) {Check prediction coordinate against fitted support};
\node[successbox,right=of coord,text width=29mm] (score) {Evaluate the observation-family posterior predictive};
\draw[arrPrimary] (post)--(cond);
\draw[arrPrimary] (cond)--(coord);
\draw[arrPrimary] (coord)--node[above,font=\scriptsize]{in support}(score);
\node[draw=Warning,fill=WarningPale,rounded corners,below=9mm of coord,text width=35mm,align=center] (err) {Outside support: reject unless an extrapolation model was specified before scoring};
\draw[arr,draw=Warning] (coord)--(err);
\node[neutralbox,below=9mm of score,text width=32mm] (match) {Boundary evaluation uses one-to-one matching; unmatched location error is NA};
\draw[arr] (score)--(match);
\end{tikzpicture}
}
\caption{Posterior prediction and evaluation. The fitted observation family and coordinate support are checked before a proper score is evaluated.}
\label{fig:posterior-prediction}
\end{figure}

\begin{definition}[Canonical one-to-one boundary matching]
\label{def:metric-f1}
Given predicted boundaries $\widehat{\mathcal B}$, references $\mathcal B^\star$, and tolerance
$\tau$, construct the bipartite graph containing edges $(\widehat b,b^\star)$ with
$|\widehat b-b^\star|\le\tau$. Select a matching with maximum cardinality; among all such matchings,
select one with minimum total absolute distance, using a fixed deterministic tie rule. Precision,
recall, and $F_1$ use the resulting matched count. No predicted or reference boundary is reused.
\end{definition}

\begin{definition}[Matched boundary-location error]
\label{def:metric-mae}
For the canonical matching $\mathcal M_\tau$, define
$\mathrm{MAE}_\tau=|\mathcal M_\tau|^{-1}\sum_{(\widehat b,b^\star)\in\mathcal M_\tau}
|\widehat b-b^\star|$ when at least one match exists. If no match exists, the quantity is
\emph{not defined} and is reported as \NA, never as zero.
\end{definition}

\begin{definition}[Held-out predictive log score]
\label{def:metric-loglik}
A held-out log score is valid only when the held-out observations are disjoint from training,
the fitted family's posterior predictive is used, and every held-out coordinate is in support or
is handled by a predeclared extrapolation policy. Failure of any condition makes the archived
number inadmissible as a predictive score.
\end{definition}

\begin{exclusionbox}[Methylation held-out computation]
The archived value $-387.50040013308154$ on $m=381$ observations is real. It is not a valid
Beta-observation held-out score because the implementation used a Gaussian predictive fallback and
silently assigned coordinates beyond the training range to the final MAP block. It is retained as
\resultid{RES-BB-RD-007Q} and excluded from predictive claims.
\end{exclusionbox}
\label{sec:prediction}

Having developed training-time posterior inference for partitions and segment parameters (Sections~\ref{sec:dp}--\ref{sec:nonconj}), we now turn to prediction for new $(X,y)$ data.
The prediction layer consumes posterior summaries (e.g., boundary posteriors, segment-parameter posteriors, and group evidences) and is therefore agnostic to whether block evidences were computed in closed form (Section~\ref{sec:families}) or via approximation (Section~\ref{sec:nonconj}). We assume we have already trained (or otherwise obtained)
a collection of $G$ group-specific BayesBreak models
$\{\mathcal{M}_g\}_{g=1}^G$, where each $\mathcal{M}_g$ comprises:
(i) a fitted segmentation posterior (e.g.\ $P_g(k\mid y)$ and boundary posteriors),
(ii) an exported MAP segmentation $\widehat{t}^{(g)}$ with segment-level posterior
hyperparameters, and (iii) optionally a Bayesian curve representation
(e.g.\ posterior mean/variance at design points).

Typical outputs are group log scores, posterior group probabilities, optional unit-level
responsibilities, exported MAP signal values, Bayes-curve values, and optional pointwise
uncertainty summaries. These prediction interfaces define the outputs consumed by the planned real-data prediction diagnostics in Section~\ref{sec:realdata}.

We pause to fix two terms that have been used informally above and that recur throughout the rest of this section.

\begin{definition}[Exported segmentation]
\label{def:exported-segmentation}
An \emph{exported segmentation} for a fitted group-specific model $\mathcal{M}_g$ is the tuple
$\big(\widehat k_g,\ \widehat t^{(g)},\ \{\widehat\theta^{(g)}_q\}_{q=1}^{\widehat k_g}\big)$,
where $\widehat k_g$ is a chosen segment count (either an estimator of $k$ such as the posterior
mode, or a deterministic choice fixed at training time), $\widehat t^{(g)}=(\widehat t^{(g)}_0,\dots,\widehat t^{(g)}_{\widehat k_g})$ is the
joint MAP boundary vector at $\widehat k_g$ produced by Algorithms~\ref{alg:map-forward}--\ref{alg:map-backtrack}, and each $\widehat\theta^{(g)}_q$ is a deterministic point summary of the
segment-$q$ posterior (typically the conjugate-family posterior mean or median of the
expectation parameter). The exported segmentation is the signal-level analogue of a fitted model;
it is what downstream prediction routines (P1 and P2) consume.
\end{definition}

\begin{definition}[Bayes curve]
\label{def:bayes-curve}
The \emph{Bayes curve} for a fitted group-specific model $\mathcal{M}_g$ evaluated at design
point $x$ is the posterior mean of the latent piecewise-constant signal under the
segment-marginalized posterior:
\[
\widehat f^{\mathrm{Bayes}}_g(x) := \mathbb{E}\!\left[\, f(x)\ \big|\ y,\, \mathcal{M}_g\,\right]
= \sum_{k=1}^{k_{\max}} P(k\mid y,\mathcal{M}_g)\,
\mathbb{E}\!\left[\, f(x)\ \big|\ y, k, \mathcal{M}_g\,\right],
\]
where the inner conditional expectation is taken over both the boundary vector $t$ and the
segment parameters $\theta_{1:k}$, and where $f(x)=\theta_{q(x)}$ is the segment-$q$ mean
parameter for the segment $q(x)$ containing $x$ under the boundary vector $t$.
\end{definition}

\begin{definition}[Segment-assignment map]
\label{def:segment-assignment-map}
For an exported boundary vector $\widehat t^{(g)}$ on ordered coordinates
$u_0^{(g)}<\cdots<u_{n_g}^{(g)}$, a segment-assignment map is a declared function
$a_g(x)\in\{1,\dots,\widehat k_g\}$ assigning each query coordinate $x$ to a training segment.
The default interval-containment rule is
$a_g(x)=q$ when $u_{\widehat t^{(g)}_{q-1}}^{(g)} < x \le u_{\widehat t^{(g)}_q}^{(g)}$, with
endpoint conventions recorded by the implementation. Queries outside the training range are either
rejected or assigned by an explicitly declared extrapolation rule; silent extrapolation is not part
of the default interface.
\end{definition}

The \emph{exported segmentation} (Definition~\ref{def:exported-segmentation}) and the
\emph{Bayes curve} (Definition~\ref{def:bayes-curve}) are distinct objects with different
operational meanings: the exported segmentation is a fixed piecewise-constant function with a
discrete change-point structure that is convenient for downstream tasks (e.g.\ region labeling,
hypothesis testing, set-valued prediction); the Bayes curve is a continuous mixture that integrates
segmentation uncertainty and is the appropriate target when smooth predictive quantities are needed.

Prediction supports two primary tasks:

\paragraph{(P1) Group membership likelihoods/posteriors.}
Given new inputs $(X,y)$, compute for each group $g$ an out-of-sample
posterior-predictive likelihood
$p(\text{new}\mid \mathcal{M}_g)$ (or its log), and convert these into posterior group
membership probabilities $P(g\mid \text{new})$ under a prior $\pi_g$.

\paragraph{(P2) Signal prediction given group.}
Given query design points $X$ and a specified (or inferred) group label $g$, return:
(i) MAP piecewise-constant signal values $\widehat{f}^{\text{MAP}}_g(X)$ including the
corresponding breakpoints and segment levels, and/or (ii) Bayesian curve values
$\widehat{f}^{\text{Bayes}}_g(X)$ with pointwise uncertainty summaries when available.

\begin{table}[H]
\centering
\small
\begin{tabular}{@{}ll@{}}
\toprule
Prediction output & Interpretation \\
\midrule
$\ell_g$ & Group-specific log posterior-predictive score. \\
$P(g\mid \mathrm{new})$ & Group posterior after combining $\ell_g$ with prior weight $\pi_g$. \\
$\widehat f^{\mathrm{MAP}}_g(X)$ & Piecewise-constant signal under the exported MAP segmentation. \\
$\widehat f^{\mathrm{Bayes}}_g(X)$ & Bayes-curve posterior mean, optionally with pointwise uncertainty. \\
Unit responsibilities & Optional probabilities for set-valued units or grouped observations. \\
\bottomrule
\end{tabular}
\caption[Prediction-layer outputs]{Prediction-layer outputs exposed by fitted BayesBreak models. These interfaces support downstream scoring and calibration separately from the empirical result tables. This is an interface-definition table; no source CSV is associated with it.}
\label{tab:prediction-outputs}
\end{table}

Crucially, we distinguish \emph{single point observations} $(x_i,y_i)$ from
\emph{multivariate / set-valued observations} where each unit contains multiple internal
pairs $(x_{u r},y_{u r})$, i.e.\ $X_u$ and $y_u$ are vectors. We also separately address
\emph{vector-valued responses} $y_i\in\mathbb{R}^d$ observed at each design point.

\subsection{Prediction inputs: pointwise vs multivariate/set-valued observations}
\label{sec:prediction-inputs}

\begin{definition}[Prediction-input classes]
\label{def:prediction-cases}
We distinguish three classes of new prediction inputs, jointly partitioning the
observation-shape design space considered in this book.
\begin{description}
\item[\textnormal{(Case A)} Pointwise observations.] A new dataset
\[
\mathcal{D}^{\mathrm{new}}_{\mathrm{pt}}
= \{(x_i^{\mathrm{new}}, y_i^{\mathrm{new}}, w_i^{\mathrm{new}})\}_{i=1}^m
\]
with $x_1^{\mathrm{new}}<\cdots<x_m^{\mathrm{new}}$, scalar (or low-fixed-dimensional)
responses $y_i^{\mathrm{new}}$, and optional exposure/precision weights
$w_i^{\mathrm{new}}$ from the observation model. Each new observation lies on a single
design point.
\item[\textnormal{(Case B)} Multivariate/set-valued observations.] A new dataset
\[
\mathcal{D}^{\mathrm{new}}_{\mathrm{mv}}
= \Big\{\big([a_u,b_u],\{(x_{ur},y_{ur},w_{ur})\}_{r=1}^{R_u}\big)\Big\}_{u=1}^U,
\]
where each \emph{unit} $u$ carries an interval support $[a_u,b_u]$ and a finite collection of
internal observations indexed by $r=1,\dots,R_u$. Prediction is performed at the unit level by
integrating or aggregating over each unit's internal observations.
\item[\textnormal{(Case C)} Vector-valued response at each design point.] A new dataset
\[
\{(x_i^{\mathrm{new}}, \mathbf{y}_i^{\mathrm{new}}, w_i^{\mathrm{new}})\}_{i=1}^m
\]
where each $\mathbf{y}_i^{\mathrm{new}}\in\mathbb{R}^d$ is a vector response (or a mixed
discrete/continuous component vector). Two sub-cases are admitted: (C-i) a factorized EF
across dimensions, and (C-ii) a joint multivariate EF with a single joint log-partition
function and conjugate or approximate block routine.
\end{description}
The three cases are mutually exclusive in the data type they consume and jointly cover the
observation shapes considered by the BayesBreak prediction interface.
\end{definition}

The remainder of this subsection treats the three cases in turn; the corresponding algorithms
appear in Section~\ref{sec:alg-prediction}.

\subsection{Posterior-predictive scoring for a fixed segment under EF conjugacy}
\label{sec:predictive-scoring}

For group membership inference, we require out-of-sample likelihoods. The most
computationally efficient and interpretable approach is to score new data against the
\emph{exported} group curve (MAP segmentation or Bayesian curve). We first derive
segment-level posterior-predictive likelihoods under EF conjugacy; these become building
blocks for both Case A and Case B.

Fix a group $g$ and consider a segment (or block) $B$ of its exported MAP segmentation.
Let the posterior for the segment parameter $\theta$ after training be the conjugate posterior
\[
p(\theta\mid \alpha^{(g)}_B,\beta^{(g)}_B)
\ \propto\
\exp\{\eta(\theta)^\top\alpha^{(g)}_B - \beta^{(g)}_B A(\theta)\},
\]
where $(\alpha^{(g)}_B,\beta^{(g)}_B)$ are the training-updated hyperparameters for that segment.
Now let $\mathcal{Y}_B^{\text{new}}$ denote the collection of new observations whose design
points fall into this segment.

\begin{proposition}[EF posterior-predictive likelihood within a fixed segment]
\label{prop:ef-predictive}
Assume the weighted EF likelihood and conjugate prior conventions of \S\ref{sec:ef} and
\S\ref{sec:notation}, and let $(S_B^{\text{new}},W_B^{\text{new}},H_B^{\text{new}})$ denote the
EF block summaries of the new data restricted to segment $B$:
\[
S_B^{\text{new}}:=\sum_{i\in B} w_i^{\text{new}} T(y_i^{\text{new}}),\qquad
W_B^{\text{new}}:=\sum_{i\in B} w_i^{\text{new}},\qquad
H_B^{\text{new}}:=\sum_{i\in B} \log h(y_i^{\text{new}},w_i^{\text{new}}).
\]
Then the posterior-predictive likelihood of the new data in segment $B$ under group $g$ is
\begin{equation}
p(\mathcal{Y}_B^{\text{new}} \mid \mathcal{M}_g)
=
\exp\{H_B^{\text{new}}\}\,
\frac{Z(\alpha_B^{(g)}+S_B^{\text{new}},\beta_B^{(g)}+W_B^{\text{new}})}
     {Z(\alpha_B^{(g)},\beta_B^{(g)})}.
\label{eq:pp-seg}
\end{equation}
\end{proposition}

\begin{proof}
The proof is identical to the single-block evidence derivation (Theorem~\ref{thm:ef-integral})
after replacing the prior hyperparameters $(\alpha_0,\beta_0)$ by the segment posterior
hyperparameters $(\alpha_B^{(g)},\beta_B^{(g)})$ learned from training data.
\end{proof}

\paragraph{Interpretation.}
Equation~\eqref{eq:pp-seg} integrates out the segment parameter $\theta$ under the
\emph{posterior} induced by training, yielding a strictly Bayesian out-of-sample likelihood
for new observations assigned to that segment. It is therefore suitable for group scoring
and group posterior computation.

\subsection{Group likelihoods for pointwise observations (Case A)}
\label{sec:group-lik-point}

\paragraph{MAP-segmentation-based scoring.}
Let $\widehat{t}^{(g)}$ denote the exported MAP boundaries for group $g$ on the relevant domain.
Assign each new design point $x_i^{\text{new}}$ to its containing MAP segment $B=B(i;g)$ and
compute segment-wise block summaries $(S_B^{\text{new}},W_B^{\text{new}},H_B^{\text{new}})$.
The group log-likelihood is then:
\begin{equation}
\ell_g^{\text{MAP}}(\mathcal{D}^{\text{new}}_{\text{pt}})
:= \log p(\mathcal{D}^{\text{new}}_{\text{pt}}\mid \mathcal{M}_g)
= \sum_{B\in \widehat{t}^{(g)}} \log p(\mathcal{Y}_B^{\text{new}}\mid \mathcal{M}_g),
\label{eq:group-lik-map}
\end{equation}
where each segment term uses \eqref{eq:pp-seg}. This yields $\mathcal{O}(m)$ evaluation time
once segment membership is determined.

\paragraph{Bayesian-curve-based scoring (optional).}
If $\mathcal{M}_g$ exports a Bayesian curve representation that provides a local posterior over
the mean parameter at each design point (or each small bin), then one can score each point
under its local posterior predictive and sum over points. This can improve robustness near
uncertain boundaries at additional computational cost.

\paragraph{Resegmentation scoring (optional, flexible).}
An alternative definition of group likelihood is to rerun the DP on the new data under
group-specific hyperparameters (treating the new dataset as another draw from group $g$'s
prior). This yields an evidence $P(\mathcal{D}^{\text{new}}_{\text{pt}}\mid g)$ that does not
condition on the trained MAP boundaries. This option is useful when groups differ primarily in
hyperparameters rather than in the learned curve geometry, but is computationally
$\mathcal{O}(k_{\max}m^2)$ per group.

\subsection{Group likelihoods for multivariate/set-valued observations (Case B)}
\label{sec:group-lik-mv}

In Case B, each unit $u$ contains multiple internal observations
$\{(x_{u r},y_{u r})\}_{r=1}^{R_u}$. Two output granularities are typically desired:
(i) \emph{unit-level} group likelihoods and group responsibilities, and (ii) a \emph{global}
group likelihood for the entire new dataset (product over units).

\paragraph{Unit-level predictive likelihood under a fixed group.}
Fix a group $g$. Partition the unit's internal design points by the MAP segments of group $g$.
For each segment $B$, compute the block summaries $(S_{u,B}^{\text{new}},W_{u,B}^{\text{new}},H_{u,B}^{\text{new}})$
restricted to those internal points. Then define the unit likelihood as:
\begin{equation}
\ell_{u g}
:= \log p(\text{unit }u \mid \mathcal{M}_g)
= \sum_{B \in \widehat{t}^{(g)}} \log p(\mathcal{Y}^{\text{new}}_{u,B}\mid \mathcal{M}_g),
\label{eq:unit-lik}
\end{equation}
where $\mathcal{Y}^{\text{new}}_{u,B}$ is the subset of unit $u$'s internal observations that
fall into segment $B$ and each segment term uses \eqref{eq:pp-seg} with summaries for that subset.

\paragraph{Dataset likelihood and independence assumption.}
\begin{assumption}[Prediction-unit conditional independence]
\label{ass:prediction-independence}
For Case B prediction, units are conditionally independent given the exported group model, and each
unit contributes only through the segment-restricted predictive factors used in \eqref{eq:unit-lik}.
\end{assumption}
Under Assumption~\ref{ass:prediction-independence}, the dataset likelihood under group $g$ is:
\[
\ell_g^{\text{MV}}(\mathcal{D}^{\text{new}}_{\text{mv}})=\sum_{u=1}^{U}\ell_{u g}.
\]
This is the direct multivariate generalization of \eqref{eq:group-lik-map} and corresponds to
the earlier ``fragment-level'' intuition, but is fully application-agnostic.

\subsection{Vector-valued responses (Case C): factorized and multivariate EF scoring}
\label{sec:prediction-multivariate-response}

Suppose each observation is $\mathbf{y}_i=(y_{i1},\dots,y_{id})$.

\paragraph{Factorized EF across dimensions (recommended default).}
Assume conditional independence across dimensions given segment parameters:
\[
p(\mathbf{y}_i\mid \boldsymbol{\theta})
=\prod_{\ell=1}^d p_\ell(y_{i\ell}\mid \theta_\ell),
\qquad
p(\boldsymbol{\theta})=\prod_{\ell=1}^d p_\ell(\theta_\ell).
\]
Then all block evidences and posterior-predictive terms factorize over $\ell$, so
log-likelihoods are sums over dimensions:
\[
\ell_g(\mathbf{y})=\sum_{\ell=1}^d \ell_{g,\ell}(y_{\cdot\ell}),
\]
where each $\ell_{g,\ell}$ is computed by \eqref{eq:group-lik-map} or \eqref{eq:unit-lik} for
dimension $\ell$. Computationally, this costs $\mathcal{O}(d)$ per observation (or per unit).

\paragraph{Multivariate EF (when available).}
If $(\mathbf{y}_i)$ admit a multivariate EF with conjugate prior on a joint parameter
$\theta\in\mathbb{R}^p$, then the same ratio-of-normalizers structure applies with vector-valued
sufficient statistics. BayesBreak's DP layer remains unchanged; only the definition of $S_{ij}$
and $Z(\alpha,\beta)$ changes.

\subsection{From group likelihoods to posterior group membership probabilities}
\label{sec:group-posterior}

\paragraph{Single-dataset membership (one group per dataset).}
If the new dataset is assumed to come from a single (unknown) group $g$, define:
\[
\ell_g := \log p(\text{new}\mid \mathcal{M}_g),
\qquad
P(g\mid \text{new})
= \frac{\pi_g\exp(\ell_g)}{\sum_{g'=1}^G \pi_{g'}\exp(\ell_{g'})}.
\]
This yields a \emph{group membership likelihood vector} (the $\ell_g$) and a
\emph{group membership posterior vector} (the $P(g\mid\text{new})$), matching the requested
prediction behavior in an application-agnostic way.

\paragraph{Unit-wise membership (one group per unit).}
If each multivariate unit $u$ is assumed to originate from one group, then define unit-wise
responsibilities:
\begin{equation}
r_{u g}
:= P(g\mid \text{unit }u)
= \frac{\pi_g\exp(\ell_{u g})}{\sum_{g'=1}^G \pi_{g'}\exp(\ell_{u g'})},
\label{eq:unit-resp}
\end{equation}
where $\ell_{u g}$ is given by \eqref{eq:unit-lik}. This is precisely the out-of-sample analog
of the E-step responsibilities in the latent-group model (Section~\ref{sec:latent-em}).

\subsection{MAP signal evaluation given group membership}
\label{sec:map-eval}

Given a group label $g$ and query design points $X^\star=\{x^\star_1,\dots,x^\star_M\}$,
BayesBreak returns either:
(i) \emph{MAP piecewise-constant predictions} or (ii) \emph{Bayesian curve predictions}.

\paragraph{MAP piecewise-constant predictions.}
Let $\widehat{t}^{(g)}$ be the exported MAP boundaries and let each MAP segment $B$ provide
the posterior mean of the mean parameter:
\[
\widehat{\mu}^{(g)}_B := \mathbb{E}[m_\star(\theta)\mid \alpha^{(g)}_B,\beta^{(g)}_B].
\]
Then the MAP signal at a query point $x^\star$ is:
\[
\widehat{f}^{\text{MAP}}_g(x^\star) = \widehat{\mu}^{(g)}_{B(x^\star;g)},
\]
where $B(x^\star;g)$ denotes the MAP segment containing $x^\star$.

\paragraph{Bayesian curve predictions.}
If the DP-derived Bayes regression curve is exported on a grid, return
$\widehat{f}^{\text{Bayes}}_g(x^\star)$ by evaluating on the nearest grid point (or by
piecewise-constant interpolation within each grid interval). When pointwise posterior
variances are available, also return pointwise credible intervals; simultaneous bands require an
additional construction and are not implied by the pointwise variances alone.

\subsection{Algorithms for prediction}
\label{sec:alg-prediction}
\label{sec:prediction-algorithms}

\begin{algorithm}[H]
\caption{PredictGroup$(\{\mathcal{M}_g\}_{g=1}^G,\ \mathcal{D}^{\text{new}},\ \pi)$}
\label{alg:predict-group}
\KwIn{
Trained group models $\{\mathcal{M}_g\}$; new data $\mathcal{D}^{\text{new}}$ (Case A or Case B);
group prior weights $\pi$ (default uniform)
}
\KwOut{
Log-likelihoods $\{\ell_g\}$ and posterior membership probabilities $\{P(g\mid \text{new})\}$;
if Case B, also unit responsibilities $\{r_{u g}\}$
}
\BlankLine
\For{$g=1$ \KwTo $G$}{
  \eIf{\textbf{Case A (pointwise)}}
  {
    Assign each $(x_i^{\text{new}},y_i^{\text{new}})$ to a MAP segment $B(i;g)$ under $\widehat{t}^{(g)}$\;
    Accumulate segment summaries $(S_B^{\text{new}},W_B^{\text{new}},H_B^{\text{new}})$\;
    Compute $\ell_g \leftarrow \sum_B \Big[ H_B^{\text{new}} + \log Z(\alpha_B^{(g)}+S_B^{\text{new}},\beta_B^{(g)}+W_B^{\text{new}})
      - \log Z(\alpha_B^{(g)},\beta_B^{(g)})\Big]$\;
  }
  {
    \tcp{Case B (multivariate/set-valued units)}
    \For{$u=1$ \KwTo $U$}{
      Partition unit $u$ internal pairs by MAP segments of group $g$; compute per-segment summaries\;
      Compute $\ell_{u g}$ via \eqref{eq:unit-lik} using \eqref{eq:pp-seg}\;
    }
    $\ell_g \leftarrow \sum_{u=1}^U \ell_{u g}$\;
  }
}
Compute $P(g\mid \text{new}) \propto \pi_g\exp(\ell_g)$ and normalize over $g$\;
\If{\textbf{Case B}}{
  Compute unit responsibilities $r_{u g}\propto \pi_g\exp(\ell_{u g})$ and normalize over $g$ for each $u$ (cf.\ \eqref{eq:unit-resp})\;
}
\end{algorithm}

\begin{algorithm}[H]
\caption{PredictMAPSignal$(\mathcal{M}_g,\ X^\star,\ \texttt{mode})$}
\label{alg:predict-map}
\KwIn{
A group model $\mathcal{M}_g$ with exported MAP segmentation $\widehat{t}^{(g)}$ and/or Bayes curve;
query points $X^\star=\{x^\star_1,\dots,x^\star_M\}$; mode $\in\{\texttt{MAP},\texttt{Bayes}\}$
}
\KwOut{
Predicted signal values at $X^\star$ (and optionally uncertainty)
}
\BlankLine
\eIf{\texttt{mode}=\texttt{MAP}}{
  \ForEach{$x^\star$ in $X^\star$}{
    Find containing MAP segment $B(x^\star;g)$ (binary search on breakpoint coordinates)\;
    Return $\widehat{f}^{\text{MAP}}_g(x^\star)=\widehat{\mu}^{(g)}_{B(x^\star;g)}$\;
  }
}{
  Evaluate/interpolate the exported Bayes regression curve at $X^\star$; optionally return posterior variances\;
}
\end{algorithm}

\subsection{Complexity}
For Case A, prediction under a fixed group is $\mathcal{O}(m)$ after segment assignment
($\mathcal{O}(m\log k)$ worst case via binary search; often $\mathcal{O}(m)$ with a streaming pointer).
For Case B, it is $\mathcal{O}(\sum_u R_u)$ per group, plus segment assignment for internal points.
Vector-valued responses add a factor of $d$ in the factorized EF construction. Importantly,
prediction does \emph{not} require rerunning the DP unless the optional resegmentation scoring
mode is enabled.
Membership posterior normalization is performed in log space by subtracting
$\texttt{logsumexp}_g(\log\pi_g+\ell_g)$; unit responsibilities use the same normalization
independently for each unit. Segment assignment is a numerical convention, not an inferential
identity: implementations must record whether binary search, streaming pointers, nearest-neighbor
assignment, or a declared extrapolation rule was used for out-of-range query points.
\subsection{Prediction diagnostics}
\label{sec:prediction-diagnostics}

Because the predictive layer reuses segment-level posteriors, the group posterior
$P(g\mid\text{new})$ derived from \eqref{eq:group-lik-map} and the pointwise predictive summaries
inherit their calibration quality from the segmentation. Use two diagnostics for real-data
deployments:
\begin{itemize}\itemsep2pt
\item \textbf{Held-out log-likelihood (HLL) trace.} Split each test sequence into blocks and
compute the per-block posterior-predictive log-likelihood under \eqref{eq:pp-seg}. The sum over
blocks is the HLL for that sequence. Compare HLL trajectories across candidate segmentations
(exported MAP vs.\ Bayes curve) and across groups; a well-calibrated model yields an HLL
ranking consistent with the $k$-posterior ranking on held-out data.
\item \textbf{PIT residuals.} For point-valued new observations with continuous predictive CDF,
evaluate the probability-integral transform $u_i=F_g(y_i^{\text{new}}\mid\mathcal{M}_g)$.
Under correct calibration the $\{u_i\}$ are Uniform$(0,1)$; histograms or Kolmogorov--Smirnov
summaries flag miscalibration (overconfident predictive widths show up as U-shaped histograms).
\end{itemize}
These diagnostics are cheap post-hoc checks. In this book they are planned outputs
for the real-data analyses in Section~\ref{sec:experiments}; the planned per-case-study results tables
identify the HLL and boundary-recovery fields that the completed analyses report.
